% Source: https://artofproblemsolving.com/community/c3h1394201_mathcounts_1996_sprint_30
% Stable identity: aops-topic-1394201-q30
Every square with vertices in this $4\times4$ array of points fits in a smallest axis-parallel bounding square with corners in the array, of some side $k\in\{1,2,3\}$. An axis-parallel square of side $k$ has exactly $k$ squares inscribed in it with a vertex on each side (the one whose vertices are its corners, and $k-1$ tilted ones, obtained by moving $j=1,\dots,k-1$ units from each corner along the sides), and each square of our kind is counted once, by its bounding square. There are $(4-k)^2$ axis-parallel squares of side $k$. So the total is
\[
\sum_{k=1}^{3}k(4-k)^2=1\cdot9+2\cdot4+3\cdot1=\boxed{20}.
\]
The figure shows a unit square (blue) and two tilted squares, with side vectors $(1,1)$ (red) and $(2,1)$ (green).

\begin{asy}
usepackage("amsmath");
size(4cm);
defaultpen(fontsize(10pt));
for(int i=1;i<=4;++i) for(int j=1;j<=4;++j) dot((i,j));
draw((-0.2,0)--(4.6,0),Arrow(TeXHead)); draw((0,-0.2)--(0,4.6),Arrow(TeXHead));
draw((1,1)--(2,1)--(2,2)--(1,2)--cycle,blue+linewidth(0.9));
draw((3,1)--(4,2)--(3,3)--(2,2)--cycle,red+linewidth(0.9));
draw((1,2)--(3,1)--(4,3)--(2,4)--cycle,heavygreen+linewidth(0.9));
label("$1$",(1,0),S);label("$4$",(4,0),S);label("$1$",(0,1),W);label("$4$",(0,4),W);
label("\mbox{}",(min(currentpicture,user=true)+max(currentpicture,user=true))/2,black); // keeps the text after the figure black when LaTeX typesets the labels inline
\end{asy}
